\documentclass[11pt]{article}
\usepackage[margin=1in]{geometry}
\usepackage{amsmath}
\usepackage{amssymb}
\usepackage{amsthm}
\usepackage{enumitem}
\usepackage{xcolor}
\usepackage{tcolorbox}

\title{Detailed Breakdown of Question 3c\\2021 ORF309 Midterm 1}
\author{}
\date{}

\begin{document}

\maketitle

\section{Problem Statement}

Find the probability that you get exposed to COVID \textbf{before} receiving your first dose of the vaccine.

\section{Key Random Variables}

\begin{tcolorbox}[colback=blue!5!white,colframe=blue!75!black,title=Random Variables]
\begin{enumerate}[label=\textbf{\arabic*.}]
    \item $X_i$: Indicator random variable for COVID exposure on day $i$
    \begin{itemize}
        \item $P(X_i = 1) = 0.01$ (exposed on day $i$)
        \item $P(X_i = 0) = 0.99$ (not exposed on day $i$)
    \end{itemize}

    \item $N$: The day you receive your first vaccine dose
    \begin{itemize}
        \item $N \sim \text{Geometric}(0.05)$
        \item Each day you have a 5\% chance of getting a vaccine slot
    \end{itemize}

    \item $A$: The event that you do \textbf{NOT} get exposed to COVID before the first dose
\end{enumerate}
\end{tcolorbox}

\section{Solution Strategy}

The solution uses the \textbf{Law of Total Probability}, conditioning on when you receive the vaccine ($N = n$).

\subsection{Why Condition on $N$?}

Since $N$ is a \emph{random variable} (you don't know in advance when you'll get the vaccine), we must consider all possible days $n = 1, 2, 3, \ldots$ and weight each by its probability.

\section{Step-by-Step Calculation}

\subsection{Step 1: Express probability for a specific day $n$}

Using the Law of Total Probability:
\begin{equation}
P(A) = \sum_{n=1}^{\infty} P\left(\sum_{i=1}^{n-1} X_i = 0 \,\Big|\, N = n\right) \cdot P(N = n)
\end{equation}

The term $\sum_{i=1}^{n-1} X_i = 0$ means you were \textbf{NOT} exposed on any day from 1 to $n-1$.

\subsection{Step 2: Use independence}

Since $X_i$ and $N$ are independent events:
\begin{equation}
P\left(\sum_{i=1}^{n-1} X_i = 0 \,\Big|\, N = n\right) = P\left(\sum_{i=1}^{n-1} X_i = 0\right) = 0.99^{n-1}
\end{equation}

This is because you need to avoid exposure for $(n-1)$ consecutive days, each with probability $0.99$.

\subsection{Step 3: Apply Geometric distribution for $N$}

For a Geometric$(p)$ random variable with parameter $p = 0.05$:
\begin{equation}
P(N = n) = 0.05 \cdot 0.95^{n-1}
\end{equation}

\subsection{Step 4: Combine and simplify}

Substituting into the equation from Step 1:
\begin{align}
P(A) &= \sum_{n=1}^{\infty} 0.99^{n-1} \cdot 0.05 \cdot 0.95^{n-1}\\
&= \sum_{n=1}^{\infty} 0.05 \cdot (0.99 \cdot 0.95)^{n-1}\\
&= 0.05 \cdot \sum_{n=1}^{\infty} (0.9405)^{n-1}
\end{align}

\subsection{Step 5: Use geometric series formula}

For $|r| < 1$, the geometric series formula gives:
\begin{equation}
\sum_{n=1}^{\infty} r^{n-1} = \frac{1}{1-r}
\end{equation}

Applying this with $r = 0.9405$:
\begin{align}
P(A) &= \frac{0.05}{1 - 0.95 \cdot 0.99}\\
&= \frac{0.05}{1 - 0.9405}\\
&= \frac{0.05}{0.0595}\\
&\approx 0.8403
\end{align}

\subsection{Final Answer}

The probability of getting exposed \textbf{before} the first dose is:
\begin{tcolorbox}[colback=green!5!white,colframe=green!75!black,title=Final Answer]
\begin{equation}
P(\text{exposed before first dose}) = 1 - P(A) = 1 - \frac{0.05}{1 - 0.95 \cdot 0.99} \approx 0.1597
\end{equation}

Or equivalently: $\boxed{1 - \frac{0.05}{1 - 0.95 \cdot 0.99}}$
\end{tcolorbox}

\section{Key Insights}

\begin{enumerate}
    \item \textbf{Why condition on $N$?} Because $N$ is random---you don't know when you'll get the vaccine, so you must account for all possible days.

    \item \textbf{Why use independence?} $X_i$ (daily COVID exposure) and $N$ (vaccine arrival) are independent events, which allows us to separate the conditional probability.

    \item \textbf{Physical interpretation}: There's roughly an 84\% chance you avoid COVID exposure before getting your first vaccine dose, which means about a 16\% chance you get exposed before being vaccinated.
\end{enumerate}

\section{Common Mistakes (from Grading Comments)}

\begin{itemize}
    \item \textbf{Treating $N$ as fixed} (loses 4.0 points): Some students fixed the date $N = n$ and computed the probability only for that specific $n$, without recognizing that $N$ itself is a random variable.

    \item \textbf{Not using Law of Total Probability}: This is essential when dealing with a random variable that affects the outcome.

    \item \textbf{Computational errors}: Mistakes in applying the geometric series formula or combining terms.
\end{itemize}

\end{document}
